\chapter{Kết luận}

Khóa luận nghiên cứu cách biểu diễn lịch sử thao tác bằng các sự kiện gắn với thực thể, bằng chứng quan sát và thời điểm khả dụng. Các mệnh đề hình thức chỉ bảo đảm quy tắc cập nhật và tính không nhìn trước của phép chiếu theo các giả định đã nêu; chúng không chứng minh độ chính xác của bộ ghi hay lợi ích điều khiển.

Về thực nghiệm, đối sánh vòng kín trên 20 trạng thái unseen ghi nhận F0 đạt 4/40 lượt thành công, F1 đạt 0/40; F1 không đạt mốc nào ở 37/40 lượt so với 30/40 của F0, còn median wall time lần lượt khoảng 1.068 s và 213 s. B0 raw bổ sung đạt 0/40 trên cùng initial-state hashes, nhưng chỉ 115/134 source hash trùng và policy runtime khác đối sánh vòng kín nên kết quả chỉ mang tính mô tả. Kiểm tra truy vấn theo sự kiện xác nhận event-triggered requery đã chạy trong closed loop: có 56 lần gọi sớm ở 22/24 episode; cả 56 log có mã evidence nhưng chỉ 41 lần thêm ID mới, còn 15 lần gắn ``History incomplete'' và giữ ID từ quan sát trước. Hai arm đều đạt 1/24 thành công; chênh lệch SCR mean là $+0{,}056$ với bootstrap 95\% CI $[-0{,}125; 0{,}250]$.

Ở nhóm hiệu chuẩn của phép kiểm tra truy vấn theo sự kiện trên states 12--19, preflight nghiêm ngặt bị chặn vì 7/14 hash nguồn writer không khớp bản kê của adapter. Sau đó hai arm chạy với \texttt{writer\_override=true}; mỗi arm hoàn tất 16 episode, đạt 1/16 thành công và khớp hash trạng thái ở cả 16 cặp. Shadow có SCR mean 0,188 và 2/16 lượt overshoot; event có SCR mean 0,083 và 0/16 overshoot. Hai arm được khởi chạy từ Git revision khác nhau, nên đây là kết quả chẩn đoán theo cấu hình override, không thay cho calibration nghiêm ngặt hay cohort xác nhận. Một lượt T1 thành công được minh họa bằng khung trích video ở Chương 4.

Phép đánh giá này đo mức duy trì hiệu năng sau tinh chỉnh trên LIBERO-10 gốc bằng đối sánh với checkpoint nền B0 trước khi tinh chỉnh trong khóa luận. Phép đánh giá này ghi nhận SmolVLA B0 37/50 và F1-off 23/50 trên LIBERO-10 local. Trên X-VLA, B0 đạt 47/50, F0 đạt 44/50 và F1-off đạt 46/50 trong 50 bộ ba ghép cặp. Kết quả X-VLA cho thấy không quan sát thấy suy giảm lớn trên roster này, nhưng cohort gần mức trần, chỉ có một seed cho mỗi trạng thái và không đủ để khẳng định tương đương. Trong đánh giá duy trì hiệu năng sau tinh chỉnh, các cấu hình nhận chỉ dẫn gốc khi suy luận. F1-off giữ trọng số F1 đã tinh chỉnh với lời nhắc bộ nhớ. Trên X-VLA, chưa có cohort rollout LIBERO-Mem đủ để kết luận về hiệu quả bộ nhớ trên mô hình này. Các kết quả đều theo giao thức nghiên cứu cục bộ, không phải điểm LIBERO chính thức và không tự xác lập quan hệ nhân quả.

Đóng góp của nghiên cứu là tách phép đo bộ ghi, nội dung memory thực sự được chiếu vào lời nhắc và kết quả hành động của policy. Kết quả hiện có cho thấy event supervisor chạy thật, nhưng chưa chứng minh memory hoặc truy vấn sớm làm tăng success; nguyên nhân các chênh lệch hiệu năng vẫn chưa được kiểm định. Các lượt vòng kín dùng SmolVLA trong mô phỏng cục bộ theo giao thức phát triển của đề tài.

Hướng tiếp theo là lặp đối sánh vòng kín và kiểm tra truy vấn theo sự kiện trên roster xác nhận khóa trước, nhiều seed và nguồn writer được khóa chặt trước preflight; đồng thời thu nhãn độc lập cho intent, target, effect và mốc commit. Việc đối chiếu bộ đếm sự kiện với bộ đánh giá hiện mới dừng ở mức mô tả, còn bộ kết xuất lời nhắc mới được kiểm tra trên trace cố định; cần nhãn mù và đánh giá tác động lên hành động trước khi kết luận về độ tin cậy hay hiệu quả. Khóa luận đã hoàn tất hai can thiệp bỏ bộ nhớ tại một truy vấn trên các trạng thái được chọn; bước tiếp theo là mở rộng roster và thêm đối chứng giả hoặc thay thế fact có cùng mức biến đổi văn bản. Cần đánh giá memory-on trên X-VLA bằng rollout LIBERO-Mem đủ cohort. Đánh giá retention cũng cần đối chứng được thiết kế trước về tỷ lệ trộn dữ liệu và định tuyến tham số.
